\section{Trajectories with a bounded bare offset}
\label{sec:trajectories}

The construction above chooses the microscopic coupling by prescribing
its value after the final blocking step.  To treat every sufficiently
large microscopic coupling, we instead allow a bounded displacement
from an explicit reference sequence.  We shall prove that all these
trajectories remain admitted.  We shall also identify which displacements
give the same effective density at every fixed layer as the microscopic
depth tends to infinity.  This second conclusion requires a summable
sensitivity estimate for the error in the kinetic increment.

\subsection{Common reference scales and canonical increments}

We keep the observations and representations of
Sections~\ref{sec:setup}--\ref{sec:rg}.  In particular, $L$ is a fixed
sufficiently large power of two, and
\[
 \gamma=\frac{\log L}{2\pi},\qquad
 H_j=H+\gamma j,\qquad
 \mathfrak m_j=\lceil(\log H_j)^2\rceil,\qquad
 p_j=(\log H_j)^{P_0},\qquad t_j=H_j^{-1/2}p_j.
\]
A depth-$N$ trajectory starts at $b_N=\beta$ and runs through layers
$N,N-1,\ldots,0$.  Its precise coupling is admitted at layer $j$ when
$b_j\in[H_j/2,2H_j]$, and strict admission means that both inequalities
are strict.  The reference scales, rather than $b_j$, determine the
geometric thresholds and masks throughout a comparison.

There are three trajectory types, denoted by $P\in\{r,+,-\}$.
Type $r$ uses only ordinary observations of side $L$.  Type $+$ or $-$
uses first the inclined observation of side $5L$ in
Section~\ref{sec:free}, with respective rotation
\[
 O_+=\frac15\begin{pmatrix}3&-4\\4&3\end{pmatrix},
 \qquad O_-=O_+^{-1},
\]
and then uses ordinary observations.  A history records the observations
already made, including their coordinate frames.  In a comparison we
identify the frames of the layer under consideration and use the same
ordinary observation in the remaining common steps.

We use the bulk formulas and their compatible finite-period
representations.  Compatibility means that every observation descends
to the corresponding quotient lattice, while admission of a period
requires its shortest nonzero translation to satisfy the lower bound
in Section~\ref{setup:scales}.  The comparisons below use common
reference scales and the common refinements and zero padding of
Section~\ref{sec:setup}.  They apply to bulk lists and simultaneously
to the finite-period lists whenever the periods being compared are
compatible.  In particular, one may fix a common terminal period and
compare the resulting common periods at all fixed bottom layers.

Recall briefly what the discrepancy measures.  A retained density has
the form \eqref{eq:setup-effective-density}: it contains the precise
kinetic coupling, a finite vector of canonical polynomial coefficient
lists, regular local errors, and a covering sum for configurations with
large bond differences.  The regular errors are bounded with eight
derivatives on their graph domains; the covering activities are bounded
by a sum over supports meeting a site.  Their caps are
\[
 \delta_j=H_j^{-2.05},\qquad w_j=\exp(-p_j^{1/4});
\]
the complete definitions are
\eqref{eq:setup-error-list-norm} and
\eqref{eq:setup-covering-norm}.  For two bulk representations, $u_j$ denotes
the positively weighted list discrepancy
\eqref{eq:setup-density-discrepancy}.  The regular-error difference
uses seven derivatives and is divided by $\delta_j$; the covering
difference is divided by $w_j$.  The canonical coefficients are
compared without their displayed powers of the precise coupling.
We put $\lambda_j=b_{2,j}-b_{1,j}$, where each precise coupling is
defined by its bulk extraction and is shared by the corresponding
compatible finite-period representations.

For a comparison that also involves finite periods, choose the common
terminal periods to be compared and follow each of them through the
layers.  Write $u_j^{\rm bulk}$ and $u_j^T$ for
\eqref{eq:setup-density-discrepancy} in the bulk and on the layer-$j$
period corresponding to terminal period $T$, respectively.  In this case we use
\begin{equation}
 u_j=\max\bigl\{u_j^{\rm bulk},\ \sup_Tu_j^T\bigr\}.
 \label{eq:offset-family-discrepancy}
\end{equation}
The supremum may be restricted to the finitely many periods in the
comparison, or taken over a compatible family obeying the common
admission bounds.  The estimates are uniform in periods, and
$\lambda_j$ and the history error are common to this whole family.
The cap bounds are preserved under this maximum; the step estimates
for this convention are justified below.  In particular a finite-period comparison always
retains control of the bulk data that determine the kinetic increment;
no norm on an isolated finite torus is used to bound a bulk extraction.
Fixed positive comparison weights give
\begin{equation}
 \|f_{2,j}-f_{1,j}\|_{7,j,A}\le C_L\delta_j u_j,
 \qquad u_j\le C_L
 \quad\hbox{for two representations obeying the common caps.}
 \label{eq:offset-discrepancy-caps}
\end{equation}
The volume scalar is omitted from $u_j$, as it cancels from normalized
expectations.

The order of parameters will matter.  Choose the contraction slack
$\upsilon>0$ small enough for Theorem~\ref{thm:rg-step} and so that
\begin{equation}
 \frac{\upsilon}{2}<\min\{1,\eta/(2\pi)\},
 \label{eq:offset-slack}
\end{equation}
where $\eta$ is the gain in Proposition~\ref{prop:preliminary}.
Choose the support and derivative requirements, $L$, the canonical
caps and comparison weights, and $P_0$ in the order specified in
Section~\ref{subsec:rg-closure}.  All logarithmic powers in the
estimates below are then fixed.  Two bounds on the bare displacement,
$R$ and $R'$, will be chosen below, before increasing $H$.
Constants denoted by $C_L$ may depend on these already fixed
representation parameters, but not on $H$, the depth, the layer,
the period, or the precise couplings in their bands.  Once $H$ has
been fixed, we allow constants denoted by $C_H$ to depend on it.

\begin{proposition}[Canonical increments and comparison estimates]
\label{prop:rg-input}
For every type $P\in\{r,+,-\}$, initialize the nearest-neighbor model
as in Section~\ref{sec:shooting}.  As long as the input couplings lie
in their admitted bands, the exact observations produce representations
with all renewed shape caps and the recursion
\begin{equation}
 b_{j-1}=b_j+\alpha^P_{N-j}
              +\frac{\kappa^P_{N-j}}{b_j}+\Delta_j,
 \qquad |\Delta_j|\le C_LH_j^{-1.05}.
 \label{eq:offset-flow}
\end{equation}
There are $C_L<\infty$, $0<\sigma_1<1$, and a number
$\kappa_\infty$ such that
\begin{equation}
 |\alpha_h^P+\gamma|+|\kappa_h^P-\kappa_\infty|
       \le C_L\sigma_1^h \qquad(h\ge0).
 \label{eq:offset-increment-convergence}
\end{equation}
These sequences are independent of the precise running couplings,
and remain fixed when $H$ is increased.  The two inclined types have
the same scalar sequences.

On common ordinary steps, with discrepancy variables defined above,
\begin{align}
 u_{j-1}&\le q u_j+C_L\operatorname{poly}(H_j)\epsilon_h
                        +e_j|\lambda_j|,\nonumber\\
 |\lambda_{j-1}-\lambda_j|
     &\le C_Lu_j+C_L\operatorname{poly}(H_j)\epsilon_h
                        +e_j|\lambda_j|,
 & e_j&=C_L(\log H_j)^C/H_j.
 \label{eq:offset-weak-comparison}
\end{align}
If the histories share their last $m$ ordinary observations,
$\epsilon_h$ may be taken to be $r_0^m$, where $r_0=A_0/L^2$.
The choices can be made with
\begin{equation}
 \max\{q,r_0\}<\rho<L^{-2+\upsilon}<1
 \label{eq:offset-contraction-width}
\end{equation}
and strict spare width.  Here $q$ and $\rho$ are fixed bounds valid
for every subsequent sufficiently large $H$.  The corresponding free kernels and their
lifts compare in the exponential kernel norms of
Proposition~\ref{prop:free-bounds} with the same history bound.
\end{proposition}

\begin{proof}
The initialization in Section~\ref{sec:shooting} has zero canonical
coefficients and zero regular errors; its covering construction obeys
the covering cap.  Theorem~\ref{thm:rg-step} renews these caps and
gives the coupling recursion whenever its input is admitted.  It
produces the output representation before the next coupling-band
condition is imposed, a fact that will permit a first-exit argument.

The prescription in Section~\ref{subsec:rg-canonical} assigns the
canonical coefficients to the finite triangular operation
$\mathbf P\mapsto\mathcal C_h(\mathbf P)$.  After the powers of
$b^{-1/2}$ have been sorted, that operation depends only on the
incoming canonical coefficients and the free history.  The remaining
difference from the exact integral is retained in the regular errors
and covering activities, with its affine correction in $\Delta$.
The orbit starting from zero therefore fixes the sequences
$\alpha_h^P,\kappa_h^P$ independently of the precise coupling and
of $H$.  Lemma~\ref{lem:rg-canonical-contraction} and the history
contraction in Proposition~\ref{prop:free-bounds} give their
exponential convergence, as in Section~\ref{sec:shooting};
Lemma~\ref{lem:kinetic-increment} identifies their common first
limit as $-\gamma$.  Reflection interchanges the two inclined
prescriptions and preserves these scalar extractions, so their
scalar sequences agree.  This does not identify their coefficient
tensors in a fixed coordinate frame.

Equation~\eqref{eq:offset-weak-comparison} follows from
\eqref{eq:orig-5} with the simultaneous bulk and period convention.
To verify this convention through the normalization, first compare
the raw lists in the bulk and on each selected period using the
dominating discrepancy \eqref{eq:offset-family-discrepancy}.
Extract the difference of the kinetic corrections from the bulk raw
lists, with the factor $Ct_{j-1}^{-2}$ in
Section~\ref{subsec:rg-closure}.  The same bulk corrections are then
used to normalize every finite-period output.  Their regular reset
costs $Ct_{j-1}^{2}$ times this difference, cancelling the extraction
factor; the other reset discrepancies are the history and coupling
terms already present in \eqref{eq:orig-5}.  Missing bulk corrections
on a given period keep their complete winding records and the spare
support exponent, as in the period argument of
Section~\ref{subsec:rg-closure}; their off-mask parts retain the
covering estimates.  These operations introduce only the fixed
normalization multipliers already allowed before choosing $L$.
Thus each output discrepancy is bounded by
$qu_j+C_L\operatorname{poly}(H_j)\epsilon_h+e_j|\lambda_j|$,
with a common contraction coefficient $q$ below $\rho$ after the
triangular weights and large parameters are chosen with spare width.
Taking their maximum proves the first family inequality;
the coupling inequality uses the bulk term directly.
The history and kernel conclusions are
Proposition~\ref{prop:free-bounds}.  The spare width in
\eqref{eq:offset-contraction-width} is the choice made in
Proposition~\ref{prop:endpoint-matching}.  The uniform $o_H(1)$
terms in the step estimates can be absorbed into a fixed upper
contraction bound $q$ strictly below $\rho$; increasing $H$ later
does not change either chosen rate.
\end{proof}

\subsection{Summable sensitivity of the kinetic error}

The coefficient $e_j$ in
\eqref{eq:offset-weak-comparison} tends to zero but is not summable
over layers.  For matching trajectories from their behavior at large
$j$, we need finer information about the noncanonical increment
$\Delta_j$.  The raw-error estimate used in the proof of the exact
step supplies it.

\begin{lemma}[Sensitivity of the extracted kinetic error]
\label{lem:sensitivity}
For two admitted compatible representation families, including their
prescribed bulk lists, on a common ordinary step with common reference
scales, the discrepancies in
Proposition~\ref{prop:rg-input} satisfy
\begin{equation}
 |\Delta_{2,j}-\Delta_{1,j}|
 \le C_L H_j^{-1.01}(u_j+|\lambda_j|)
             +C_L\operatorname{poly}(H_j)\epsilon_h.
 \label{eq:offset-sensitivity}
\end{equation}
The constant is uniform over the bulk prescriptions and their
compatible admitted finite-period representations, once $H$ is
sufficiently large.
\end{lemma}

\begin{proof}
Before the constant and affine terms of the regular remainder are
extracted, let $R_j^\Delta$ denote the layer-$(j-1)$, seven-derivative
norm of the difference of its two bulk raw lists.  The Taylor comparison proved in
Section~\ref{subsec:rg-closure} gives the estimate
\cite[Equation~(5.39)]{OpenAI-O4}, with the improved coefficient
$q_1=C/L^2+o_H(1)$ established there for $S^2$:
\begin{equation}
 R_j^\Delta\le q_1\delta_j^\Delta
   +C_L\operatorname{poly}(\mathfrak m_j,p_j)\delta_j
                      (\epsilon_h+\lambda_g)
   +C_Lg^{4.5}(u_j+\epsilon_h+\lambda_g).
 \label{eq:offset-raw-discrepancy}
\end{equation}
Here $\delta_j^\Delta$ is the unnormalized bulk input regular-error
discrepancy; $g$ may be taken comparable to either $b_{i,j}^{-1/2}$;
and $\lambda_g\le C_L|\lambda_j|/H_j$.  By
\eqref{eq:offset-discrepancy-caps},
$\delta_j^\Delta\le C_LH_j^{-2.05}u_j$, while
$g\le C_LH_j^{-1/2}$ in the common bands.

The affine coefficient is the linear second-derivative functional
\eqref{eq:setup-affine-coefficient}, evaluated in the bulk.  Once
the coordinate frames and reference domains have been identified,
this is the same functional on both raw lists; it does not change
with the free kinetic kernel.  The extraction estimate in
Section~\ref{subsec:rg-closure} therefore gives
\begin{equation}
 |\Delta_{2,j}-\Delta_{1,j}|
       \le C t_{j-1}^{-2}R_j^\Delta
       \le C_L H_j R_j^\Delta.
 \label{eq:offset-affine-extraction}
\end{equation}
There is no support-cardinality loss in this estimate: a unit affine
direction is bounded by the polynomially weighted directions in
\eqref{eq:setup-error-directions}, so its two derivative slots cost
only $t_{j-1}^{-2}$.  The seven derivatives in the discrepancy norm
are more than sufficient.  The final reset of the canonical
prefactors has zero affine Hessian, as verified in
Section~\ref{subsec:rg-closure}, and hence contributes no additional
kinetic extraction.  Bulk extraction also explains the uniformity
in the period: missing or folded corrections remain in the winding
lists instead of changing $\Delta_j$.  Equation~\eqref{eq:offset-affine-extraction}
uses the bulk raw-list discrepancy, which is bounded using
$u_j^{\rm bulk}\le u_j$ from
\eqref{eq:offset-family-discrepancy}.  The inclusion of this bulk
term is essential when some labels wind or disappear on a finite period.

Multiplying \eqref{eq:offset-raw-discrepancy} by $C_LH_j$ bounds its
nonhistory terms by
\[
 C_L(H_j^{-1.05}+H_j^{-1.25})u_j
 +C_L\bigl((\log H_j)^C H_j^{-2.05}
                         +H_j^{-2.25}\bigr)|\lambda_j|.
\]
Any remaining fixed logarithmic powers are absorbed by the strict
power margins when $H$ is increased.  In particular this is bounded
by $C_LH_j^{-1.01}(u_j+|\lambda_j|)$.  The history terms can be
bounded by $C_L\operatorname{poly}(H_j)\epsilon_h$, proving
\eqref{eq:offset-sensitivity}.  The argument uses the raw-error
comparison for $S^2$ from Section~\ref{subsec:rg-closure}, rather
than an application of an $O(4)$ step theorem to the present model.
\end{proof}

\subsection{Admission and retuning}

The exponentially decaying canonical transients define the finite
constants
\begin{equation}
 A_P=\sum_{h=0}^{\infty}(\alpha_h^P+\gamma),\qquad
 A_{\rm abs}=\max_{P\in\{r,+,-\}}
                  \sum_{h=0}^{\infty}|\alpha_h^P+\gamma|,
 \qquad d=A_+-A_r=A_--A_r.
 \label{eq:offset-transients}
\end{equation}
These constants are fixed before $H$ is increased.  Recall the
reference sequence from Section~\ref{sec:shooting}:
\begin{equation}
 \vartheta_j=H_j-\frac{\kappa_\infty}{\gamma}\log(H_j/H),
 \qquad
 \vartheta_j-\vartheta_{j-1}
       =\gamma-\frac{\kappa_\infty}{H_j}+O_L(H_j^{-2}).
 \label{eq:offset-profile}
\end{equation}
The bare offset of a depth-$N$ trajectory is
$s=\beta-\vartheta_N$.  Choose once and for all
\begin{equation}
 R>4(|d|+\gamma+1),\qquad
 R'>R+2A_{\rm abs}+2,\qquad I=(-R,R).
 \label{eq:offset-range}
\end{equation}
We shall use $|s|\le R$ for the limiting comparisons and the larger
range $|s|\le R'$ for retuning.  Both bounds have now been chosen
independently of $H$.

\begin{proposition}[Uniform admission for bounded offsets]
\label{prop:admission}
For $H$ sufficiently large, every depth $N\ge1$, every type
$P\in\{r,+,-\}$, and every bare coupling
$\beta=\vartheta_N+s$ with $|s|\le R'$ give a strictly admitted
trajectory.  Uniformly over these choices,
\begin{equation}
 b_j=\vartheta_j+D_j,\qquad
 \max_{0\le j\le N}|D_j|\le C_{L,R'}.
 \label{eq:offset-admission}
\end{equation}
Moreover,
\begin{align}
 D_{j-1}-D_j&=\alpha^P_{N-j}+\gamma+E_{j,N},\nonumber\\
 |E_{j,N}|&\le C_L\left[
 H_j^{-1.05}
 +H_j^{-2}\bigl(\log(2+H_j/H)+|D_j|\bigr)
 +H_j^{-1}\sigma_1^{N-j}\right],
 \label{eq:offset-error-recursion}
\end{align}
and
\begin{equation}
 \sup_{N,P,\,|s|\le R'}\sum_{j=1}^{N}|E_{j,N}|
       =o_H(1).
 \label{eq:offset-error-sum}
\end{equation}
The constants in \eqref{eq:offset-admission} stay bounded as $H$
is increased, with all earlier parameters fixed.
\end{proposition}

\begin{proof}
The starting coupling is strictly in band for $H$ large, uniformly
in $N$ and $|s|\le R'$.  Indeed, for $x\ge1$ the ratio
$\log x/x$ is bounded, and hence
$|\vartheta_N-H_N|/H_N\le C_L/H$.
Consider a trajectory down to a possible first output outside its
band.  Every step forming that output has an admitted input, so
\eqref{eq:offset-flow} applies.  On such an input,
\[
 \left|\frac1{b_j}-\frac1{H_j}\right|
 \le C_LH_j^{-2}
          \bigl(\log(2+H_j/H)+|D_j|\bigr).
\]
Subtract \eqref{eq:offset-profile} from the coupling recursion and
use \eqref{eq:offset-increment-convergence}.  This gives
\eqref{eq:offset-error-recursion}, including the step forming a
possible first exit.

The sums required for this estimate are uniform in depth:
\begin{align*}
 \sum_{j\ge1}H_j^{-1.05}&\le C_LH^{-0.05},&
 \sum_{j\ge1}H_j^{-2}&\le C_LH^{-1},\\
 \sum_{j\ge1}H_j^{-2}\log(2+H_j/H)&\le C_LH^{-1},&
 \sum_{j=1}^{N}H_j^{-1}\sigma_1^{N-j}&\le C_LH^{-1}.
\end{align*}
The first three follow by integral comparison for
$H_j=H+\gamma j$; the last follows from the geometric series and
$H_j\ge H$.  Let $M_D$ be the maximum of $|D_j|$ on the segment,
including its last output.  Summing from $D_N=s$ and using
\eqref{eq:offset-transients} gives
\[
 M_D\le R'+A_{\rm abs}+o_H(1)+C_LH^{-1}M_D.
\]
Take $H$ large enough to absorb the final term.  This bounds $M_D$
by a constant depending only on $L,R'$ and the previously fixed
parameters.  Since also
\[
 \sup_{j\ge0}\frac{|\vartheta_j-H_j|+C_{L,R'}}{H_j}
       \longrightarrow0\qquad(H\longrightarrow\infty),
\]
every output covered by this bound lies strictly in its band.
There can be no first exit, which proves admission and
\eqref{eq:offset-admission}.  Substitution of this uniform bound
for $|D_j|$ in \eqref{eq:offset-error-recursion}, followed by the
same four summations, proves \eqref{eq:offset-error-sum}.
\end{proof}

\begin{lemma}[Retuning at an arbitrary depth]
\label{lem:retuning}
After increasing $H$ once more, let an ordinary trajectory have any
depth $N\ge1$ and bare offset $|s|\le R$.  For every integer $K\ge1$
there is $s_K\in[-R',R']$ such that the ordinary trajectory of
depth $K$ with bare coupling $\vartheta_K+s_K$ has exactly the
same terminal coupling $b_0$.
\end{lemma}

\begin{proof}
For every trajectory in Proposition~\ref{prop:admission}, summing
\eqref{eq:offset-error-recursion} to zero gives
\begin{equation}
 \left|b_0-H-s\right|\le A_{\rm abs}+o_H(1).
 \label{eq:offset-terminal-range}
\end{equation}
Let $T_K(t)$ be the terminal coupling of the ordinary depth-$K$
trajectory started at $\vartheta_K+t$.  Proposition~\ref{prop:admission}
gives strict admission for the whole interval $[-R',R']$.
The continuity argument in the proof of Proposition~\ref{prop:shooting}
therefore makes $T_K$ continuous on this interval: the reference
thresholds are fixed, and the exact integrations and representation
maps are continuous in the precise coupling on each finite admitted
run.  Take $H$ so that the uniform $o_H(1)$ in
\eqref{eq:offset-terminal-range} has absolute value less than one.
By \eqref{eq:offset-range},
\begin{align*}
 T_K(-R')&\le H-R'+A_{\rm abs}+1
                  <H-R-A_{\rm abs}-1\le b_0,\\
 T_K(R')&\ge H+R'-A_{\rm abs}-1
                  >H+R+A_{\rm abs}+1\ge b_0.
\end{align*}
The intermediate value theorem supplies $s_K$.  No monotonicity or
uniqueness of the shooting map is needed.
\end{proof}

\subsection{Matching at each fixed layer}

Retuning gives an exact equality of terminal couplings at finite
depths.  A different comparison identifies limits without such a
retuning.  Its hypothesis compensates the transient displacement
$A_P$ associated with the choice of initial observation.

\begin{proposition}[Fixed-layer matching]
\label{prop:matching}
Fix $H$ sufficiently large as above.  For $i\in\{1,2\}$, let
$P_i\in\{r,+,-\}$ be fixed and let a sequence of depth-$N_{i,n}$
trajectories have bare couplings
$\vartheta_{N_{i,n}}+s_{i,n}$, where
$N_{i,n}\longrightarrow\infty$ and $|s_{i,n}|\le R'$.
Suppose
\begin{equation}
 (s_{2,n}+A_{P_2})-(s_{1,n}+A_{P_1})\longrightarrow0.
 \label{eq:offset-matching-hypothesis}
\end{equation}
Compare their bulk representations in common layer coordinates,
using the same ordinary observations on their common bottom steps.
Then, for every fixed integer $j\ge0$,
\begin{equation}
 u_{j,n}\longrightarrow0,\qquad
 \lambda_{j,n}\longrightarrow0.
 \label{eq:offset-fixed-layer-matching}
\end{equation}
The same statement holds on each fixed common compatible terminal
period and its corresponding bottom-layer periods.
\end{proposition}

\begin{proof}
All the trajectories are admitted by Proposition~\ref{prop:admission}.
At layer $j$, the histories share at least
$\min(N_{1,n},N_{2,n})-1-j$ most recent ordinary observations,
once the depths are sufficiently large.  Thus their history error
tends to zero at each fixed layer.  The normalized caps and
\eqref{eq:offset-admission} also give constants, independent of
$j,n$, bounding $u_{j,n}$ and $|\lambda_{j,n}|$ whenever the layer
is present.  Define
\[
 U_j=\limsup_{n\to\infty}u_{j,n},\qquad
 W_j=\limsup_{n\to\infty}|\lambda_{j,n}|.
\]
Both sequences are uniformly bounded.

We first establish the boundary condition for $W_j$ at large
layers.  The recursion in Proposition~\ref{prop:admission} gives
\[
 D_{i,j}=s_{i,n}
       +\sum_{h=0}^{N_{i,n}-j-1}(\alpha_h^{P_i}+\gamma)
       +\sum_{m=j+1}^{N_{i,n}}E_{i,m,N_{i,n}}.
\]
For fixed $j$ the transient sum tends to $A_{P_i}$ as
$n\to\infty$.  The error bound, with
\eqref{eq:offset-admission} substituted, gives
\begin{equation}
 W_j\le C_{L,R'}\sum_{m>j}
       \left[H_m^{-1.05}
          +H_m^{-2}\bigl(\log(2+H_m/H)+1\bigr)\right]
       +\frac{C_L}{H_j}.
 \label{eq:offset-boundary-at-infinity}
\end{equation}
Here \eqref{eq:offset-matching-hypothesis} cancels the two limiting
offsets, and the final term bounds each transient sum by
\[
 \sum_{m=j+1}^{N}H_m^{-1}\sigma_1^{N-m}
                 \le \frac1{H_j(1-\sigma_1)}.
\]
The right side of \eqref{eq:offset-boundary-at-infinity} tends to
zero as $j\to\infty$, with $H$ fixed.  Consequently
\begin{equation}
 W_j\longrightarrow0\qquad(j\longrightarrow\infty).
 \label{eq:offset-W-boundary}
\end{equation}
This argument does not require the complete error sum at fixed $H$
to vanish as the depth increases.

At each fixed $j\ge1$, take upper limits in the first inequality of
\eqref{eq:offset-weak-comparison}.  The history forcing disappears.
For the coupling, subtract \eqref{eq:offset-flow} for the two
trajectories and use Lemma~\ref{lem:sensitivity}.  At this fixed
layer the two $\alpha$ values tend to $-\gamma$ and the two
$\kappa$ values to $\kappa_\infty$.  The remaining denominator
variation is bounded by $C_LH_j^{-2}|\lambda_{j,n}|$.
It follows that
\begin{equation}
 U_{j-1}\le qU_j+e_jW_j,\qquad
 W_{j-1}\le W_j+a_j(U_j+W_j),\qquad
 a_j=C_LH_j^{-1.01}.
 \label{eq:offset-limsup-recursions}
\end{equation}

Put $U=\sup_{j\ge0}U_j$, $W=\sup_{j\ge0}W_j$, and
$e_* =\sup_{j\ge1}e_j$.  Iterating the first inequality through
$m$ layers above $k$ yields
\[
 U_k\le q^m U_{k+m}
             +W\sum_{r=1}^{m}q^{r-1}e_{k+r}.
\]
Boundedness and $q<1$ allow $m\to\infty$, giving
$U\le e_*W/(1-q)$.  Next sum the second inequality of
\eqref{eq:offset-limsup-recursions} from $k+1$ to $k+m$.
Letting $m\to\infty$ and using \eqref{eq:offset-W-boundary}
gives
\[
 W_k\le\sum_{j>k}a_j(U_j+W_j),\qquad
 W\le\left(\sum_{j\ge1}a_j\right)(U+W).
\]
No limit has been interchanged with an infinite sum: the limiting
recurrences were first obtained at fixed layers and then iterated.
Finally,
\[
 e_*\longrightarrow0,\qquad
 \sum_{j\ge1}a_j\le C_LH^{-0.01}\longrightarrow0
                  \qquad(H\longrightarrow\infty).
\]
We choose $H$, once for all the trajectories, so that
\[
 \left(\sum_{j\ge1}a_j\right)
                  \left(1+\frac{e_*}{1-q}\right)<1.
\]
The two supremum inequalities force $W=0$ and then $U=0$,
which proves \eqref{eq:offset-fixed-layer-matching}.  The proof
uses the bulk norms and the simultaneous finite-period norms from
Proposition~\ref{prop:rg-input}; it therefore applies to each fixed
compatible terminal period as stated.
\end{proof}

For later use, all further requirements that $H$ be large, including
the local source and terminal-alignment estimates, are imposed before
any depth or volume thresholds are selected.  The bounds $R,R'$,
the transient constants, and the reference prescriptions remain fixed
throughout these choices.
